\chapter{Background}
\label{chap:background}

This chapter summarizes the parts of the typesetting pipeline that matter for
incremental rebuilds.

\section{The \TeX{} Pipeline}
\label{sec:pipeline}

\TeX{} turns a stream of characters into tokens, expands macros, and builds
lists of boxes and glue. Paragraphs are broken into lines by an algorithm
that minimizes a global cost over the whole paragraph
\citep{knuth1981breaking}. Figure~\ref{fig:pipeline} shows the stages.

\begin{figure}[ht]
  \centering
  \setlength{\unitlength}{1mm}
\begin{picture}(140,22)
  \put(0,4){\framebox(26,14){\strut Source}}
  \put(26,11){\vector(1,0){12}}
  \put(38,4){\framebox(26,14){\strut Tokens}}
  \put(64,11){\vector(1,0){12}}
  \put(76,4){\framebox(26,14){\strut Boxes}}
  \put(102,11){\vector(1,0){12}}
  \put(114,4){\framebox(26,14){\strut Pages}}
\end{picture}

  \caption{The stages of a typesetting run. An incremental system tries to
           resume from a saved state inside the engine instead of restarting
           at the source.}
  \label{fig:pipeline}
\end{figure}

\subsection{Auxiliary Files}
\label{sec:aux}

Cross-references and citations are resolved through auxiliary files that one
run writes and the next run reads. This is why a fresh document needs more
than one pass, and it is the main source of dependencies that point backwards
in the document \citep{mittelbach2004companion}.

\section{Documents as Programs}
\label{sec:programs}

A \LaTeX{} source file is a program whose output is a document. The idea of
treating a document and its code as a single text is older than \LaTeX{}
itself \citep{knuth1984literate}. It also means that, in general, the effect
of an edit cannot be predicted without running the program.

\section{Measuring Information in an Edit}
\label{sec:information}

A small edit carries little information relative to the full source. If an
edit selects one of $N$ equally likely alternatives, its information content
is
\begin{equation}
  I = \log_2 N
  \label{eq:information}
\end{equation}
bits \citep{shannon1948mathematical}. Equation~\ref{eq:information} gives a
lower bound on the work a rebuild must do; a batch run does far more.
